
\subsubsection*{2.5.1}
\phantomsection
\label{sga3:VIIB:2.5.1}

Suppose again that \(k\) is an arbitrary pseudocompact ring. Since
the functor
\[
  X\longmapsto X_e
\]
of \SGARefLink{sga3:VIIB:1.6.F}{1.6.F} commutes with finite products, it sends every formal group
\(G\) to an étale formal group \(G_e\). The morphism
\(p:X\to X_e\) of that result is functorial in \(X\), so in this
case
\[
  p:G\longrightarrow G_e
\]
is a morphism of formal groups.

Consider its kernel \(\ker(p)\).\footnote{%
\label{sga3:VIIB:nde:97}The editors supplied the details of the
original argument that follow.} It is the fiber product of the
diagram
\begin{center}
  \phantomsection
  \label{sga3:VIIB:2.5.1:kernel-fiber-square}
  \begin{tikzcd}[row sep=2.5em,column sep=3.8em]
{} & G \arrow[d,"p"] \\
\operatorname{Spf}(k) \arrow[r,"\varepsilon"]
  & G_e
\end{tikzcd}

\end{center}
where \(\varepsilon\) is the identity section of \(G_e\). Since
\(p\) induces a bijection on the underlying sets, it follows
\textup{(cf.\ \SGARefLink{sga3:VIIB:1.2}{1.2})} that the underlying set of
\(\ker(p)\) is the image of \(\operatorname{Spf}(k)\) under
\(\varepsilon\), and that, for every point \(s\) of
\(\operatorname{Spf}(k)\), the local algebra of \(\ker(p)\) at
\(\varepsilon(s)\) is
\[
  \mathcal O_{G,\varepsilon(s)}
  \widehat\otimes_{\mathcal O_{G_e,\varepsilon(s)}}
  \kappa(s).
\]
At each point \(\varepsilon(s)\), the residue field of
\(\mathcal O_{G,\varepsilon(s)}\) is \(\kappa(s)\), and hence
\[
  \mathcal O_{G_e,\varepsilon(s)}=\kappa(s),
  \qquad
  \mathcal O_{\ker(p),\varepsilon(s)}
  =
  \mathcal O_{G,\varepsilon(s)}.
\]

For these reasons, we call \(\ker(p)\) the \emph{infinitesimal
neighborhood of the identity} of \(G\), and write
\[
  \ker(p)=G_0.
\]
We have thus obtained the exact sequence
\begin{center}
  \phantomsection
  \label{sga3:VIIB:2.5.1:infinitesimal-exact-sequence}
  \(
1\longrightarrow G_0\xrightarrow{\ \mathrm{incl.}\ }G
\xrightarrow{\ p\ }G_e.
\)

\end{center}

In what follows, we say that \(G\) is \emph{infinitesimal} if
\[
  G=G_0.
\]
\footnote{\label{sga3:VIIB:nde:98}The editors added the sentence that
follows.}
This is equivalent to saying that, for every
\(s\in\operatorname{Spf}(k)\), or equivalently for every continuous
morphism \(k\to\kappa\), where \(\kappa\) is a field with the
discrete topology, the identity is the unique element of
\(G(\kappa(s))\), respectively of \(G(\kappa)\).

Suppose in addition that \(G\) is topologically flat over
\(k\).\footnote{\label{sga3:VIIB:nde:99}The editors supplied the
details that follow.} By \SGARefLink{sga3:VIIB:1.6.1}{1.6.1}, the morphism
\[
  p:G\longrightarrow G_e
\]
is topologically flat. Since it is bijective, \SGARefLink{sga3:VIIB:1.3.1}{Proposition 1.3.1}
shows that it is an effective epimorphism. Consequently, \(G_e\) is
identified with the quotient \(G/G_0\). We have proved:

\medskip
\noindent\textbf{Corollary.}
\phantomsection
\label{sga3:VIIB:2.5.1:corollary}
{\itshape Let \(G\) be a formal group topologically flat over \(k\).
Then \(G_e\) is identified with the quotient \(G/G_0\); in other
words, there is an exact sequence of formal groups}
\begin{center}
  \phantomsection
  \label{sga3:VIIB:2.5.1:short-exact-sequence}
  \(
1\longrightarrow G_0\xrightarrow{\ \mathrm{incl.}\ }G
\xrightarrow{\ p\ }G_e\longrightarrow1.
\)

\end{center}
